\section{Thực trạng}

Trong những năm gần đây, sự phát triển vượt bậc của các mô hình ngôn ngữ lớn (Large Language Models – LLMs) đã tạo ra bước chuyển biến quan trọng trong lĩnh vực trí tuệ nhân tạo và xử lý ngôn ngữ tự nhiên. Các mô hình ngôn ngữ lớn như GPT, Gemini hay Claude có khả năng hiểu ngữ cảnh, tổng hợp thông tin và tạo sinh văn bản với độ tự nhiên cao, cho phép xây dựng các hệ thống hỏi đáp và trợ lý tri thức có khả năng tương tác linh hoạt với người dùng. Sự ra đời của các mô hình AI tiên tiến này đang thúc đẩy một làn sóng mới trong cách chúng ta thu thập và tìm kiếm thông tin, giúp cho việc khám phá một miền tri thức mới trở nên đơn giản và dễ dàng hơn bao giờ hết.

Nhờ sự thúc đẩy của AI, việc số hóa, xây dựng hạ tầng tri thức số ngày càng được quan tâm, và lĩnh vực di sản văn hóa truyền thống đang là ưu tiên hàng đầu trong công cuộc chuyển đổi số đó. Nghệ thuật Việt Nam vốn là kết tinh của lịch sử, văn hóa và đời sống tinh thần của dân tộc, được hình thành và phát triển qua nhiều thế kỷ với những đặc trưng riêng biệt. Và trong tất cả những loại hình nghệ thuật âm nhạc dân gian lâu đời ấy, nghệ thuật Chèo nổi lên như một hình thức nghệ thuật đặc sắc, chứa đựng cả một hệ thống tri thức phong phú liên quan đến vở diễn, nhân vật, nghệ sĩ, phong cách biểu diễn và bối cảnh lịch sử, ... Tuy nhiên, phần lớn tri thức này hiện nay vẫn đang tồn tại dưới dạng văn bản phi cấu trúc, phân tán trong nhiều nguồn tài liệu khác nhau. Việc truy vấn và tổng hợp thông tin thường đòi hỏi người nghiên cứu phải đọc và đối chiếu nhiều tài liệu, gây tốn thời gian và khó đảm bảo tính toàn diện.

Hầu hết các mô hình AI đều được huấn luyện sẵn trên các bộ dữ liệu lớn mang tính phổ quát cao, đồng thời cũng không có cơ chế truy cập trực tiếp tới cơ sở tri thức bên ngoài tại thời điểm suy luận. Chính điều đó tạo ra một trong những hạn chế đáng chú ý của LLMs - hiện tượng sinh thông tin không có căn cứ (hallucination). Chính vì hạn chế này, khi người dùng muốn tìm kiếm và khám phá về nghệ thuật Chèo thông qua AI, những rủi ro về tính chính xác và độ tin cậy của thông tin có thể xuất hiện, bởi nghệ thuật Chèo vốn chưa được chuyển đổi số một các toàn diện, các thông tin về Chèo vốn dĩ không xuất hiện trong các bộ dữ liệu huấn luyện trước đó của AI.

Trong bối cảnh đó, phương pháp sinh ngôn ngữ có tăng cường truy xuất (Retrieval-Augmented Generation - RAG) được đề xuất nhằm bổ sung tri thức bên ngoài vào quá trình tạo sinh của LLMs. Tuy nhiên, phần lớn các hệ thống RAG hiện nay vẫn dựa trên truy xuất đoạn văn bản dựa trên mã hóa, mà chưa khai thác triệt để cấu trúc quan hệ trong dữ liệu. Điều này dẫn đến hạn chế khi xử lý các truy vấn đa bước hoặc các truy vấn yêu cầu tổng hợp thông tin từ nhiều thực thể có liên kết logic chặt chẽ.

Do đó, việc tích hợp đồ thị tri thức (Knowledge Graph – KG), một hình thức biểu diễn tri thức có cấu trúc dựa trên thực thể và quan hệ – vào kiến trúc RAG trở thành một hướng nghiên cứu tiềm năng. Knowledge Graph cho phép biểu diễn tri thức dưới dạng đồ thị, trong đó các thực thể và mối quan hệ được xác định rõ ràng, hỗ trợ suy luận đa bước và truy vết nguồn dữ liệu. Tuy nhiên, việc kết hợp Knowledge Graph với LLMs một cách hiệu quả vẫn là một thách thức nghiên cứu đang được quan tâm.

\section{Bài toán thực tế}

Từ thực trạng nêu trên, bài toán đặt ra là phải xây dựng một đồ thị tri thức của miền nghệ thuật Chèo, chuẩn hóa nó và tạo ra một hệ thống có khả năng khai thác tri thức, sử dụng mô hình ngôn ngữ lớn để tạo sinh câu trả lời tự nhiên, chính xác và có thể truy vết. Bài toán không chỉ đơn thuần là xây dựng và truy xuất thông tin, mà còn yêu cầu kết nối các thực thể và quan hệ trong Knowledge Graph nhằm hỗ trợ suy luận đa bước.

Ví dụ, với một truy vấn ngôn ngữ tự nhiên $q$ từ người dùng, hệ thống cần thực hiện các bước sau: (i) xác định các thực thể liên quan trong Knowledge Graph; (ii) truy xuất tập hợp các quan hệ và thực thể lân cận cần thiết để trả lời câu hỏi; (iii) tổ chức ngữ cảnh có cấu trúc từ dữ liệu đồ thị; và (iv) sử dụng mô hình ngôn ngữ lớn để tạo sinh câu trả lời dựa trên ngữ cảnh đã được kiểm soát. Bài toán này đòi hỏi sự phối hợp giữa biểu diễn tri thức có cấu trúc và khả năng suy luận ngôn ngữ tự nhiên.

Thách thức chính của bài toán nằm ở việc cân bằng giữa độ bao phủ tri thức và khả năng kiểm soát. Nếu truy xuất quá ít thông tin, câu trả lời có thể thiếu chính xác; nếu truy xuất quá nhiều, ngữ cảnh trở nên dài và khó xử lý đối với LLM. Đồng thời, hệ thống cũng cần phải xác định được khả năng kết hợp thông tin truy vấn được để có thể trả lời các câu hỏi của người dùng nhanh và chính xác nhất. Ngoài ra, hệ thống cần đảm bảo rằng mọi thông tin trong câu trả lời đều có thể truy vết về Knowledge Graph, nhằm giảm thiểu hiện tượng hallucination.

Do đó, khóa luận này tập trung vào việc xây dựng một đồ đồ thị tri thức của nghệ thuật Chèo và thiết kế một kiến trúc GraphRAG, trong đó Knowledge Graph đóng vai trò trung tâm trong quá trình truy xuất và tổ chức tri thức, còn LLMs đóng vai trò diễn giải và tổng hợp.

\section{Đóng góp của khóa luận}

Các đóng góp chính của khóa luận này được thể hiện qua ba nội dung cốt lõi sau:

Thứ nhất, khóa luận xây dựng và chuẩn hóa một Knowledge Graph chuyên biệt cho nghệ thuật Chèo, trong đó các thực thể và quan hệ được tổ chức một cách có hệ thống nhằm phản ánh cấu trúc tri thức đặc thù của miền. Quá trình chuẩn hóa này không chỉ đảm bảo tính nhất quán của dữ liệu mà còn tạo nền tảng cho các nghiên cứu tiếp theo trong lĩnh vực số hóa và khai thác tri thức văn hóa truyền thống.

Thứ hai, khóa luận đề xuất một kiến trúc GraphRAG tích hợp cơ chế truy xuất đa độ hạt cùng với phương pháp tổ chức ngữ cảnh có kiểm soát. Kiến trúc này cho phép khai thác hiệu quả cấu trúc quan hệ của đồ thị nhằm hỗ trợ suy luận đa bước, đồng thời duy trì năng lực sinh ngôn ngữ tự nhiên của các mô hình ngôn ngữ lớn mà không yêu cầu tái huấn luyện.

Thứ ba, khóa luận thiết kế một bộ dữ liệu đánh giá (benchmark) chuyên biệt cho miền nghệ thuật Chèo, bao gồm việc phân loại truy vấn theo mức độ phức tạp và xây dựng ground truth với cấu trúc chi tiết. Bộ benchmark này tạo điều kiện cho việc đánh giá định lượng một cách hệ thống, đồng thời hỗ trợ phân tích sâu hiệu năng của mô hình.

\section{Cấu trúc khóa luận}

Khóa luận được tổ chức thành năm chương chính, bên cạnh các phần mở đầu và tài liệu tham khảo, nhằm đảm bảo trình bày một cách hệ thống từ cơ sở lý thuyết đến phương pháp đề xuất và đánh giá thực nghiệm. Các chương tiếp theo được tổ chức như sau:

\begin{itemize}
    \item Chương 2 – Trình bày cơ sở lý thuyết về nghệ thuật Chèo, các hình thức biểu diễn tri thức và các phương pháp RAG/GraphRAG, đồng thời tổng quan các công trình liên quan.
    \item Chương 3 – Đề xuất kiến trúc hệ thống GraphRAG, bao gồm xây dựng Knowledge Graph, cơ chế truy xuất, tạo sinh và thiết kế bộ dữ liệu đánh giá.
    \item Chương 4 – Triển khai hệ thống, thiết kế thực nghiệm và đánh giá hiệu năng thông qua các độ đo và so sánh với phương pháp truyền thống.
    \item Chương 5 – Tổng kết kết quả đạt được, nêu đóng góp, hạn chế và định hướng phát triển trong tương lai.
\end{itemize}

